% =============================================================================
% APPENDIX GGG: METHOD-CONFIG: Model Configurator
% =============================================================================
% Module: BCM2_METHOD_CONFIG
% Category: METHOD
% Language: English
% Version: 1.0 (January 2026)
% Status: Draft
%
% This appendix defines the Model Configurator---a systematic tool for generating
% behavioral models from configuration dimensions. The Configurator uses mapping
% tables to auto-populate model parameters based on Level, Geography, Domain,
% and Scope Type selections.
%
% =============================================================================

% =============================================================================
% CROSS-REFERENCE STRUCTURE
% =============================================================================
%
% CHAPTER LINKAGE (Required):
% - Primary Chapter: Chapter 4: From Framework to Model
% - Secondary Chapters: Chapter 10 (Applications), Chapter 9 (Context)
%
% DEPENDENCIES (what this appendix REQUIRES):
% - Appendix UNMAPPED_DSN (METHOD-DESIGN): Workflow that uses configurator
% - Appendix UNMAPPED_REG (METHOD-REGISTRY): Stores generated models
% - Appendix UNMAPPED_WHO (CORE-WHO): Level definitions
% - Appendix UNMAPPED_HOW (CORE-HOW): Complementarity
% - Appendix UNMAPPED_WAT (CORE-WHAT): Utility dimensions
% - Appendix CTW (CORE-WHEN): Context dimensions
% - Appendix UNMAPPED_AWA (CORE-AWARE): Awareness function [NEW]
% - Appendix UNMAPPED_REA (CORE-READY): Willingness function [NEW]
% - Appendix UNMAPPED_STA (CORE-STAGE): Behavioral change journey [NEW]
%
% DEPENDENTS (what REQUIRES this appendix):
% - Appendix UNMAPPED_DSN (METHOD-DESIGN): Uses configurator for model generation
% - All DOMAIN- appendices: Domain-specific mappings
%
% =============================================================================

\begin{tcolorbox}[colback=purple!5!white, colframe=purple!75!black,
    title=Chapter Linkage]

\textbf{Primary Chapter:} Chapter 4: From Framework to Model
\begin{itemize}[nosep]
\item Section 4.3: Model Configuration $\leftrightarrow$ Section \ref{sec:ggg-dimensions}
\item Section 4.4: Parameter Mapping $\leftrightarrow$ Section \ref{sec:ggg-mappings}
\end{itemize}

\textbf{Secondary Chapters:}
\begin{itemize}[nosep]
\item Chapter 9: Context --- provides $\Psi$ dimension structure
\item Chapter 10: Applications --- uses configurator for domain models
\item Chapter 11: Awareness --- provides $A(\cdot)$ function (AWARE)
\item Chapter 12: Willingness --- provides $WAX$, $\theta$ (READY)
\item Chapter 13: Behavioral Change Journey --- provides stage dynamics (STAGE)
\end{itemize}

\textbf{Reading Guide:}
\begin{itemize}[nosep]
\item For workflow context: Read Appendix UNMAPPED_DSN first
\item For configuration logic: This appendix provides complete specification
\item For model storage: See Appendix UNMAPPED_REG (METHOD-REGISTRY)
\end{itemize}

\end{tcolorbox}

\chapter{METHOD-CONFIG: Model Configurator}
\label{app:method-config}

\begin{abstract}
This appendix defines the Model Configurator---a systematic tool for generating
\textbf{10C-complete} behavioral models from four configuration dimensions: Level ($L$),
Geography ($\Psi_{geo}$), Domain, and Scope Type. The Configurator uses seven mapping
tables to auto-populate all 10C parameters: context ($\Psi$), utility dimensions ($C$),
complementarity ($\gamma$), functional form, awareness ($A$), willingness ($WAX$, $\theta$),
and stage dynamics ($\tau$, $\rho$, $\kappa$). This transforms model design from
an expert-dependent art into a structured, reproducible process with complete 10C coverage.
\end{abstract}


\begin{tcolorbox}[
    colback=orange!5!white,
    colframe=orange!75!black,
    title={\textbf{Appendix Scope: Ziel / In-Scope / Out-of-Scope / Constraints}},
    fonttitle=\bfseries
]

\textbf{Ziel (Objective):}
\begin{itemize}[nosep]
    \item Provide systematic analysis for METHOD integration
\end{itemize}

\textbf{In-Scope:}
\begin{itemize}[nosep]
    \item Core analysis and findings
    \item Framework integration points
\end{itemize}

\textbf{Out-of-Scope:}
\begin{itemize}[nosep]
    \item Detailed empirical replication $\rightarrow$ METHOD appendices
\end{itemize}

\textbf{Constraints:}
\begin{itemize}[nosep]
    \item Analysis bounded by available data
\end{itemize}

\end{tcolorbox}

\begin{tcolorbox}[colback=blue!5!white, colframe=blue!75!black,
    title=Quick Reference: Model Configurator (10C Complete)]
\textbf{Appendix:} CFG (METHOD)


\textbf{Core Question:} How do we systematically generate \textbf{10C-complete} model specifications from high-level choices?

\textbf{Key Principle:}
\begin{equation}
\text{Config}(L, \Psi_{geo}, \text{Domain}, \text{Scope}) \rightarrow M(\Psi, C, \gamma, f, A, WAX, \theta, \tau, \rho, \kappa)
\end{equation}

\textbf{Key Concepts:}
\begin{itemize}[nosep]
\item \textbf{4 Configuration Dimensions}: Level, Geography, Domain, Scope Type
\item \textbf{7 Mapping Tables}: Dimension combinations $\rightarrow$ 10C model parameters
\item \textbf{10C Auto-Population}: WHO, WHAT, HOW, WHEN, WHERE, AWARE, READY, STAGE
\item \textbf{3+1 Architecture}: Each dimension offers 3 options + custom
\end{itemize}

\textbf{Cross-References:} EEE (Workflow), FFF (Registry), AAA (WHO), B (HOW), C (WHAT), V (WHEN), AU (AWARE), AV (READY), AW (STAGE)
\end{tcolorbox}


% =============================================================================
%                    PART 2: CORE CONTENT
% =============================================================================

%%%%%%%%%%%%%%%%%%%%%%%%%%%%%%%%%%%%%%%%%%%%%%%%%%%%%%%%%%%%%%%%%%%%%%%%%%%%%%%
\section{The Fundamental Question}
\label{sec:ggg-fundamental}
%%%%%%%%%%%%%%%%%%%%%%%%%%%%%%%%%%%%%%%%%%%%%%%%%%%%%%%%%%%%%%%%%%%%%%%%%%%%%%%

\subsection{Why This Matters}

The EEE Workflow defines \textit{how} to design a model step-by-step. But at each
step, the modeler must make choices: Which $\Psi$ dimensions? Which $C$ components?
Which functional form? These choices require expertise and can vary widely between
practitioners.

The Configurator standardizes these choices through \textbf{mapping tables} that
encode organizational knowledge about which parameters work for which situations.

\subsection{The Core Problem}

\begin{tcolorbox}[colback=red!5!white, colframe=red!75!black,
    title=The Fundamental Question]
\textbf{How do we translate high-level configuration choices (Level, Geography,
Domain, Scope) into concrete model specifications ($\Psi$, $C$, $\gamma$, functional form)?}
\end{tcolorbox}


%%%%%%%%%%%%%%%%%%%%%%%%%%%%%%%%%%%%%%%%%%%%%%%%%%%%%%%%%%%%%%%%%%%%%%%%%%%%%%%
\section{The Four Configuration Dimensions}
\label{sec:ggg-dimensions}
%%%%%%%%%%%%%%%%%%%%%%%%%%%%%%%%%%%%%%%%%%%%%%%%%%%%%%%%%%%%%%%%%%%%%%%%%%%%%%%

\subsection{Dimension 1: Level ($L$)}

The welfare hierarchy level from CORE-WHO (Appendix UNMAPPED_WHO).

\begin{tcolorbox}[colback=white, colframe=black,
    title=Dimension 1: Level Selection]
\textbf{3+1 Options:}
\begin{enumerate}
\item[$\square$] \textbf{L=0 Individual:} Single person decisions (enrollment, purchase, adherence)
\item[$\square$] \textbf{L=1 Household:} Family/household decisions (savings, insurance, housing)
\item[$\square$] \textbf{L=3 Organization:} Firm/institution decisions (adoption, change, investment)
\item[$\square$] \textbf{Custom:} L=2 (Team), L=5 (Market), L=6 (Nation), other
\end{enumerate}

\textbf{Implication:} Level determines which utility dimensions are relevant and
how cross-level effects ($\gamma^{cross}$) operate.
\end{tcolorbox}

\subsection{Dimension 2: Geography ($\Psi_{geo}$)}

The geographic/regulatory context.

\begin{tcolorbox}[colback=white, colframe=black,
    title=Dimension 2: Geography Selection]
\textbf{3+1 Options:}
\begin{enumerate}
\item[$\square$] \textbf{Switzerland (CH):} BVG/FINMA regulatory, high digital adoption, high institutional trust
\item[$\square$] \textbf{Austria (AT):} ASVG/FMA regulatory, medium state orientation, EU-aligned
\item[$\square$] \textbf{Germany (DE):} BetrAVG/BaFin regulatory, high state orientation, complex bureaucracy
\item[$\square$] \textbf{Custom:} Other DACH region, EU, global
\end{enumerate}

\textbf{Implication:} Geography sets regulatory $\Psi$, cultural $\Psi$, and
institutional trust parameters.
\end{tcolorbox}

\subsection{Dimension 3: Domain}

The application domain determines which utility dimensions are salient.

\begin{tcolorbox}[colback=white, colframe=black,
    title=Dimension 3: Domain Selection]
\textbf{3+1 Options:}
\begin{enumerate}
\item[$\square$] \textbf{Vorsorge (Pension/Retirement):} F, E, S dominant; long-term temporal
\item[$\square$] \textbf{Gesundheit (Health):} P, E dominant; present bias relevant
\item[$\square$] \textbf{Finance (Banking/Investment):} F dominant; risk parameters critical
\item[$\square$] \textbf{Custom:} HR/Org, Public, Energy, Retail, other
\end{enumerate}

\textbf{Implication:} Domain determines $C$ dimension weights and typical $\gamma$ interactions.
\end{tcolorbox}

\subsection{Dimension 4: Scope Type}

The behavioral outcome type determines functional form.

\begin{tcolorbox}[colback=white, colframe=black,
    title=Dimension 4: Scope Type Selection]
\textbf{3+1 Options:}
\begin{enumerate}
\item[$\square$] \textbf{Decision:} Binary choice (yes/no, enroll/not, buy/not) → Logistic
\item[$\square$] \textbf{Continuous:} Amount, frequency, intensity → Linear/CES
\item[$\square$] \textbf{Process:} Journey, adoption stages, change phases → Dynamic/BCJ
\item[$\square$] \textbf{Custom:} Hybrid, multi-stage, other
\end{enumerate}

\textbf{Implication:} Scope Type determines functional form and estimation approach.
\end{tcolorbox}


%%%%%%%%%%%%%%%%%%%%%%%%%%%%%%%%%%%%%%%%%%%%%%%%%%%%%%%%%%%%%%%%%%%%%%%%%%%%%%%
\section{Mapping Tables}
\label{sec:ggg-mappings}
%%%%%%%%%%%%%%%%%%%%%%%%%%%%%%%%%%%%%%%%%%%%%%%%%%%%%%%%%%%%%%%%%%%%%%%%%%%%%%%

\subsection{Mapping 1: Level × Geography → $\Psi$ Defaults}

\begin{table}[htbp]
\centering
\caption{$\Psi$ Parameter Defaults by Level and Geography}
\label{tab:ggg-psi-mapping}
\renewcommand{\arraystretch}{1.2}
\small
\begin{tabular}{@{}ll|ccc@{}}
\toprule
& & \multicolumn{3}{c}{\textbf{Geography}} \\
\textbf{Level} & \textbf{$\Psi$ Dimension} & 🇨🇭 CH & 🇦🇹 AT & 🇩🇪 DE \\
\midrule
\multirow{4}{*}{L=0 Individual}
& $\Psi_{trust}$ (institutional) & 0.85 & 0.70 & 0.75 \\
& $\Psi_{digital}$ (adoption) & 0.90 & 0.80 & 0.70 \\
& $\Psi_{autonomy}$ (self-reliance) & 0.85 & 0.65 & 0.60 \\
& $\Psi_{risk}$ (tolerance) & 0.60 & 0.50 & 0.45 \\
\addlinespace
\multirow{4}{*}{L=1 Household}
& $\Psi_{trust}$ & 0.80 & 0.65 & 0.70 \\
& $\Psi_{digital}$ & 0.85 & 0.75 & 0.65 \\
& $\Psi_{social}$ (family orientation) & 0.70 & 0.80 & 0.75 \\
& $\Psi_{planning}$ (horizon) & 0.75 & 0.65 & 0.70 \\
\addlinespace
\multirow{4}{*}{L=3 Organization}
& $\Psi_{regulatory}$ (compliance) & 0.90 & 0.85 & 0.95 \\
& $\Psi_{innovation}$ (openness) & 0.80 & 0.65 & 0.60 \\
& $\Psi_{hierarchy}$ (structure) & 0.50 & 0.70 & 0.80 \\
& $\Psi_{consensus}$ (decision style) & 0.85 & 0.70 & 0.65 \\
\bottomrule
\end{tabular}
\end{table}

\subsection{Mapping 2: Domain × Level → $C$ Dimensions}

\begin{table}[htbp]
\centering
\caption{Utility Dimension Weights by Domain and Level}
\label{tab:ggg-c-mapping}
\renewcommand{\arraystretch}{1.2}
\small
\begin{tabular}{@{}ll|cccccc@{}}
\toprule
& & \multicolumn{6}{c}{\textbf{C Dimensions (FEPSDE)}} \\
\textbf{Domain} & \textbf{Level} & F & E & P & S & D & X \\
\midrule
\multirow{2}{*}{Vorsorge}
& L=0 Individual & 0.35 & 0.25 & 0.10 & 0.20 & 0.05 & 0.05 \\
& L=1 Household & 0.40 & 0.20 & 0.05 & 0.30 & 0.03 & 0.02 \\
\addlinespace
\multirow{2}{*}{Gesundheit}
& L=0 Individual & 0.15 & 0.30 & 0.35 & 0.15 & 0.03 & 0.02 \\
& L=1 Household & 0.20 & 0.25 & 0.30 & 0.20 & 0.03 & 0.02 \\
\addlinespace
\multirow{2}{*}{Finance}
& L=0 Individual & 0.50 & 0.20 & 0.05 & 0.15 & 0.08 & 0.02 \\
& L=3 Organization & 0.55 & 0.10 & 0.02 & 0.20 & 0.10 & 0.03 \\
\addlinespace
\multirow{2}{*}{HR/Org}
& L=0 Individual & 0.25 & 0.30 & 0.10 & 0.25 & 0.08 & 0.02 \\
& L=3 Organization & 0.30 & 0.15 & 0.05 & 0.35 & 0.12 & 0.03 \\
\addlinespace
\multirow{2}{*}{Public}
& L=0 Individual & 0.30 & 0.20 & 0.10 & 0.30 & 0.05 & 0.05 \\
& L=6 Nation & 0.35 & 0.10 & 0.15 & 0.25 & 0.10 & 0.05 \\
\addlinespace
\multirow{2}{*}{Energy}
& L=0 Individual & 0.35 & 0.20 & 0.15 & 0.15 & 0.05 & 0.10 \\
& L=1 Household & 0.40 & 0.15 & 0.15 & 0.15 & 0.05 & 0.10 \\
\bottomrule
\end{tabular}
\end{table}

\textbf{Legend:} F=Financial, E=Emotional, P=Physical, S=Social, D=Developmental, X=Existential

\subsection{Mapping 3: Domain → $\gamma$ Interactions}

\begin{table}[htbp]
\centering
\caption{Default Complementarity Interactions by Domain}
\label{tab:ggg-gamma-mapping}
\renewcommand{\arraystretch}{1.3}
\begin{tabular}{@{}lp{5cm}l@{}}
\toprule
\textbf{Domain} & \textbf{Primary $\gamma$ Interactions} & \textbf{Typical Values} \\
\midrule
Vorsorge & $\gamma_{F \times E}$ (security-emotion) & 0.20--0.40 \\
& $\gamma_{temporal}$ (present-future) & 0.15--0.30 \\
& $\gamma_{F \times S}$ (wealth-status) & 0.10--0.25 \\
\addlinespace
Gesundheit & $\gamma_{P \times E}$ (health-wellbeing) & 0.30--0.50 \\
& $\gamma_{temporal}$ (present bias) & 0.25--0.45 \\
& $\gamma_{P \times S}$ (health-social) & 0.10--0.20 \\
\addlinespace
Finance & $\gamma_{F \times E}$ (wealth-anxiety) & 0.15--0.35 \\
& $\gamma_{risk}$ (gain-loss asymmetry) & 0.20--0.40 \\
& $\gamma_{F \times S}$ (wealth-status) & 0.15--0.30 \\
\addlinespace
HR/Org & $\gamma_{S \times E}$ (belonging-satisfaction) & 0.25--0.45 \\
& $\gamma_{D \times F}$ (growth-compensation) & 0.15--0.30 \\
& $\gamma_{cross}$ (individual-org) & 0.10--0.25 \\
\addlinespace
Public & $\gamma_{S \times F}$ (compliance-cost) & 0.15--0.30 \\
& $\gamma_{trust}$ (institution-behavior) & 0.20--0.40 \\
\addlinespace
Energy & $\gamma_{F \times X}$ (cost-environment) & 0.10--0.30 \\
& $\gamma_{S \times X}$ (norms-environment) & 0.15--0.35 \\
& $\gamma_{temporal}$ (now-future) & 0.20--0.40 \\
\bottomrule
\end{tabular}
\end{table}

\subsection{Mapping 4: Scope Type → Functional Form}

\begin{table}[htbp]
\centering
\caption{Functional Form by Scope Type}
\label{tab:ggg-form-mapping}
\renewcommand{\arraystretch}{1.4}
\begin{tabular}{@{}llp{5cm}@{}}
\toprule
\textbf{Scope Type} & \textbf{Functional Form} & \textbf{Specification} \\
\midrule
Decision & Logistic & $P(Y=1) = \frac{1}{1+\exp(-(\beta_0 + \sum_d \beta_d C_d + \gamma \cdot \text{interactions}))}$ \\
\addlinespace
Continuous & Linear/CES & $Y = \beta_0 + \sum_d \beta_d C_d + \gamma \cdot \text{interactions} + \varepsilon$ \\
& & or CES: $Y = [\sum_d \alpha_d C_d^\rho]^{1/\rho}$ \\
\addlinespace
Process & Dynamic/BCJ & $\frac{dS}{dt} = \frac{1}{\tau}[S^*(A, WAX, \Psi) - S(t)] + \varepsilon$ \\
& & with stages $\varphi \in \{1,...,5\}$ \\
\bottomrule
\end{tabular}
\end{table}


\subsection{Mapping 5: Domain → Awareness Defaults ($A$)}

Derived from Appendix UNMAPPED_AWA (CORE-AWARE) and Chapter 11.

\begin{table}[htbp]
\centering
\caption{Awareness Parameter Defaults by Domain}
\label{tab:ggg-awareness-mapping}
\renewcommand{\arraystretch}{1.3}
\begin{tabular}{@{}lp{4cm}ll@{}}
\toprule
\textbf{Domain} & \textbf{Awareness Characteristics} & \textbf{$A_{base}$} & \textbf{Tag} \\
\midrule
Vorsorge & Low salience until retirement nears; abstract future consequences & 0.35 & [LLM] \\
Gesundheit & Varies by symptom visibility; acute > chronic & 0.50 & [EMP] \\
Finance & Depends on financial literacy; often overconfident & 0.45 & [EMP] \\
HR/Org & Workplace feedback provides signals; moderate visibility & 0.55 & [LLM] \\
Public & Policy rarely salient; requires media activation & 0.40 & [LLM] \\
Energy & Consumption invisible; delayed feedback (monthly bill) & 0.30 & [EMP] \\
\bottomrule
\end{tabular}
\end{table}

\textbf{Modulation:} $A_{eff} = A_{base} \cdot \Psi_{salience} \cdot \sigma_{trigger}$ where $\sigma_{trigger} \in \{0,1\}$ indicates contextual activation (AWX-A2a).


\subsection{Mapping 6: Domain → Willingness Defaults ($WAX$, $\theta$)}

Derived from Appendix UNMAPPED_REA (CORE-READY) and Chapter 12.

\begin{table}[htbp]
\centering
\caption{Willingness and Threshold Defaults by Domain}
\label{tab:ggg-willingness-mapping}
\renewcommand{\arraystretch}{1.3}
\begin{tabular}{@{}lllp{4cm}l@{}}
\toprule
\textbf{Domain} & \textbf{$WAX_{base}$} & \textbf{$\theta_{base}$} & \textbf{Key Barriers} & \textbf{Tag} \\
\midrule
Vorsorge & 0.40 & 0.60 & Complexity, temporal distance, inertia & [EMP] \\
Gesundheit & 0.55 & 0.50 & Pain avoidance, habit strength, social support & [EMP] \\
Finance & 0.45 & 0.55 & Loss aversion, trust, cognitive load & [EMP] \\
HR/Org & 0.50 & 0.45 & Status quo bias, political risk, effort & [LLM] \\
Public & 0.35 & 0.65 & Free-rider incentives, diffuse benefits & [LLM] \\
Energy & 0.45 & 0.50 & Comfort trade-offs, upfront costs, hassle & [EMP] \\
\bottomrule
\end{tabular}
\end{table}

\textbf{Action Condition:} $\text{Action} \Leftrightarrow WAX \geq \theta$ with gap $\Delta = WAX - \theta$ determining action probability.

\textbf{Intervention Implication:}
\begin{itemize}[nosep]
\item If $\Delta < 0$: Either increase $WAX$ (motivation) or reduce $\theta$ (friction)
\item Reducing $\theta$ often more effective than increasing $WAX$ (cf. Appendix UNMAPPED_REA, WAX-A15)
\end{itemize}


\subsection{Mapping 7: Domain × Complexity → Stage Defaults}

Derived from Appendix UNMAPPED_STA (CORE-STAGE) and Chapter 13.

\begin{table}[htbp]
\centering
\caption{Behavioral Change Journey Parameters by Domain}
\label{tab:ggg-stage-mapping}
\renewcommand{\arraystretch}{1.3}
\begin{tabular}{@{}llllll@{}}
\toprule
\textbf{Domain} & \textbf{$C_{behav}$} & \textbf{$\tau$} & \textbf{$\rho$} & \textbf{$\kappa_{S3}$} & \textbf{Tag} \\
 & (Complexity) & (Transition) & (Relapse) & (Contemplation) & \\
\midrule
Vorsorge & 0.60 & 0.08/Mo & 0.25 & 0.75 & [LLM] \\
Gesundheit & 0.65 & 0.12/Mo & 0.45 & 0.70 & [EMP] \\
Finance & 0.50 & 0.10/Mo & 0.30 & 0.65 & [LLM] \\
HR/Org & 0.55 & 0.15/Mo & 0.35 & 0.60 & [LLM] \\
Public & 0.40 & 0.06/Mo & 0.40 & 0.70 & [LLM] \\
Energy & 0.45 & 0.18/Mo & 0.35 & 0.55 & [EMP] \\
\bottomrule
\end{tabular}
\end{table}

\textbf{Legend:}
\begin{itemize}[nosep]
\item $C_{behav}$: Behavior complexity (0=trivial, 1=extreme); cf. BCJ-B3
\item $\tau$: Transition rate (monthly); higher = faster stage progression
\item $\rho$: Relapse probability; higher = more likely to regress
\item $\kappa_{S3}$: Stage 3 persistence (``contemplation trap''); higher = more stuck
\end{itemize}

\textbf{Stage Dynamics:} $\frac{dS}{dt} = \tau \cdot (WAX - \theta) \cdot (1-\kappa) + \beta \cdot \text{Social}$

\textbf{Key Insight:} Energy has highest $\tau$ (fast adoption possible) but moderate $\rho$ (relapse common without habit formation).


%%%%%%%%%%%%%%%%%%%%%%%%%%%%%%%%%%%%%%%%%%%%%%%%%%%%%%%%%%%%%%%%%%%%%%%%%%%%%%%
\section{The Configuration Algorithm}
\label{sec:ggg-algorithm}
%%%%%%%%%%%%%%%%%%%%%%%%%%%%%%%%%%%%%%%%%%%%%%%%%%%%%%%%%%%%%%%%%%%%%%%%%%%%%%%

\begin{tcolorbox}[colback=gray!5!white, colframe=gray!75!black,
    title=Algorithm CONFIG-1: Model Generation (10C Complete)]
\textbf{Input:} User selections $(L, \Psi_{geo}, \text{Domain}, \text{ScopeType})$

\textbf{Process:}
\begin{enumerate}
\item \textbf{Validate inputs:} Check all four dimensions are specified
\item \textbf{Lookup $\Psi$:} Use Table \ref{tab:ggg-psi-mapping} with $(L, \Psi_{geo})$ \hfill [WHEN]
\item \textbf{Lookup $C$:} Use Table \ref{tab:ggg-c-mapping} with $(\text{Domain}, L)$ \hfill [WHAT]
\item \textbf{Lookup $\gamma$:} Use Table \ref{tab:ggg-gamma-mapping} with $(\text{Domain})$ \hfill [HOW]
\item \textbf{Select Form:} Use Table \ref{tab:ggg-form-mapping} with $(\text{ScopeType})$
\item \textbf{Lookup $A$:} Use Table \ref{tab:ggg-awareness-mapping} with $(\text{Domain})$ \hfill [AWARE]
\item \textbf{Lookup $WAX$, $\theta$:} Use Table \ref{tab:ggg-willingness-mapping} with $(\text{Domain})$ \hfill [READY]
\item \textbf{Lookup Stage params:} Use Table \ref{tab:ggg-stage-mapping} with $(\text{Domain})$ \hfill [STAGE]
\item \textbf{Generate Model:} Assemble $M = (\Psi, C, \gamma, f, A, WAX, \theta, \tau, \rho, \kappa)$
\item \textbf{Mark sources:} Tag all parameters as ``CONFIG-derived''
\end{enumerate}

\textbf{Output:} Complete 10C model specification ready for refinement

\textbf{10C Coverage:}
\begin{itemize}[nosep]
\item WHO: $L$ (input) $\rightarrow$ Appendix UNMAPPED_WHO
\item WHAT: $C$ (Mapping 2) $\rightarrow$ Appendix UNMAPPED_WAT
\item HOW: $\gamma$ (Mapping 3) $\rightarrow$ Appendix UNMAPPED_HOW
\item WHEN: $\Psi$ (Mapping 1) $\rightarrow$ Appendix CTW
\item WHERE: All tables $\rightarrow$ Appendix UNMAPPED_EST
\item AWARE: $A$ (Mapping 5) $\rightarrow$ Appendix UNMAPPED_AWA
\item READY: $WAX$, $\theta$ (Mapping 6) $\rightarrow$ Appendix UNMAPPED_REA
\item STAGE: $\tau$, $\rho$, $\kappa$ (Mapping 7) $\rightarrow$ Appendix UNMAPPED_STA
\end{itemize}

\textbf{Post-Condition:} Model can be refined in EEE workflow, then stored in FFF registry
\end{tcolorbox}

\subsection{Configuration Flow Diagram}

\begin{figure}[htbp]
\centering
\begin{tikzpicture}[node distance=1.5cm, auto,
    input/.style={draw, rectangle, rounded corners, fill=blue!10, minimum width=3cm},
    mapping/.style={draw, rectangle, fill=yellow!10, minimum width=3cm},
    output/.style={draw, rectangle, rounded corners, fill=green!10, minimum width=3cm}]

    % Inputs
    \node[input] (level) {Level $L$};
    \node[input, right=of level] (geo) {Geography};
    \node[input, right=of geo] (domain) {Domain};
    \node[input, right=of domain] (scope) {Scope Type};

    % Mappings
    \node[mapping, below=1.5cm of level] (psi) {$\Psi$ Mapping};
    \node[mapping, below=1.5cm of geo] (c) {$C$ Mapping};
    \node[mapping, below=1.5cm of domain] (gamma) {$\gamma$ Mapping};
    \node[mapping, below=1.5cm of scope] (form) {Form Mapping};

    % Arrows to mappings
    \draw[->] (level) -- (psi);
    \draw[->] (geo) -- (psi);
    \draw[->] (level) -- (c);
    \draw[->] (domain) -- (c);
    \draw[->] (domain) -- (gamma);
    \draw[->] (scope) -- (form);

    % Output
    \node[output, below=3.5cm of geo, xshift=1.5cm] (model) {Generated Model $M(\Psi, C, \gamma, f)$};

    % Arrows to output
    \draw[->] (psi) -- (model);
    \draw[->] (c) -- (model);
    \draw[->] (gamma) -- (model);
    \draw[->] (form) -- (model);

\end{tikzpicture}
\caption{Configuration Flow: From Dimensions to Model}
\label{fig:ggg-flow}
\end{figure}


%%%%%%%%%%%%%%%%%%%%%%%%%%%%%%%%%%%%%%%%%%%%%%%%%%%%%%%%%%%%%%%%%%%%%%%%%%%%%%%
\section{DACH Reference Configurations}
\label{sec:ggg-dach}
%%%%%%%%%%%%%%%%%%%%%%%%%%%%%%%%%%%%%%%%%%%%%%%%%%%%%%%%%%%%%%%%%%%%%%%%%%%%%%%

Eight pre-validated configurations for common DACH use cases.

\begin{table}[htbp]
\centering
\caption{DACH Reference Configurations}
\label{tab:ggg-dach-configs}
\renewcommand{\arraystretch}{1.3}
\small
\begin{tabular}{@{}llllll@{}}
\toprule
\textbf{ID} & \textbf{Level} & \textbf{Geo} & \textbf{Domain} & \textbf{Scope} & \textbf{Use Case} \\
\midrule
DACH-01 & L=0 & 🇨🇭 & Vorsorge & Decision & PK Enrollment \\
DACH-02 & L=0 & 🇨🇭 & Vorsorge & Continuous & Beitragshöhe \\
DACH-03 & L=0 & 🇦🇹 & Gesundheit & Decision & Vorsorgeuntersuchung \\
DACH-04 & L=1 & 🇩🇪 & Finance & Decision & Riester Familie \\
DACH-05 & L=3 & 🇨🇭 & HR/Org & Process & Change Adoption \\
DACH-06 & L=3 & 🇩🇪 & HR/Org & Process & Digitalisierung \\
DACH-07 & L=0 & 🇦🇹 & Public & Decision & Steuer Compliance \\
DACH-08 & L=1 & 🇨🇭 & Energy & Continuous & Verbrauchsreduktion \\
\bottomrule
\end{tabular}
\end{table}

\subsection{Example: DACH-01 Configuration Output}

\begin{tcolorbox}[colback=white, colframe=black,
    title=DACH-01: PK Enrollment Switzerland (L=0, CH, Vorsorge, Decision)]
\textbf{Generated $\Psi$:}
\begin{verbatim}
psi_trust: 0.85
psi_digital: 0.90
psi_autonomy: 0.85
psi_risk: 0.60
psi_regulatory: "BVG"
\end{verbatim}

\textbf{Generated $C$ weights:}
\begin{verbatim}
C_F: 0.35  (Financial security)
C_E: 0.25  (Emotional peace of mind)
C_P: 0.10  (Physical - low relevance)
C_S: 0.20  (Social norms, employer)
C_D: 0.05  (Developmental)
C_X: 0.05  (Existential - legacy)
\end{verbatim}

\textbf{Generated $\gamma$ interactions:}
\begin{verbatim}
gamma_FxE: 0.30  (security-emotion synergy)
gamma_temporal: 0.20  (present-future tradeoff)
gamma_FxS: 0.15  (wealth-status link)
\end{verbatim}

\textbf{Generated functional form:}
\begin{equation}
P(\text{enroll}) = \frac{1}{1+\exp(-(\beta_0 + 0.35 F + 0.25 E + 0.20 S + 0.30 F \cdot E))}
\end{equation}

\textbf{Source Tag:} CONFIG-derived, requires validation
\end{tcolorbox}


%%%%%%%%%%%%%%%%%%%%%%%%%%%%%%%%%%%%%%%%%%%%%%%%%%%%%%%%%%%%%%%%%%%%%%%%%%%%%%%

%%%%%%%%%%%%%%%%%%%%%%%%%%%%%%%%%%%%%%%%%%%%%%%%%%%%%%%%%%%%%%%%%%%%%%%%%%%%%%%
\subsection{Core Axioms}
\label{sec:cfg-axioms}
%%%%%%%%%%%%%%%%%%%%%%%%%%%%%%%%%%%%%%%%%%%%%%%%%%%%%%%%%%%%%%%%%%%%%%%%%%%%%%%

\begin{tcolorbox}[colback=yellow!5!white,colframe=yellow!75!black,title=Axiom UNMAPPED_CFG-1: Fundamental Principle]
\textbf{Axiom UNMAPPED_CFG-1 (Core Assumption):}
The fundamental principle underlying this appendix is that behavioral outcomes
emerge from the interaction of individual characteristics and contextual factors.
\end{tcolorbox}

\begin{tcolorbox}[colback=yellow!5!white,colframe=yellow!75!black,title=Axiom UNMAPPED_CFG-2: Complementarity]
\textbf{Axiom UNMAPPED_CFG-2 (Complementarity):}
Components of the framework exhibit complementarity ($\gamma > 0$), meaning
that combined effects exceed the sum of individual effects.
\end{tcolorbox}



%%%%%%%%%%%%%%%%%%%%%%%%%%%%%%%%%%%%%%%%%%%%%%%%%%%%%%%%%%%%%%%%%%%%%%%%%%%%%%%
\section{Key Results}
\label{sec:cfg-results}
%%%%%%%%%%%%%%%%%%%%%%%%%%%%%%%%%%%%%%%%%%%%%%%%%%%%%%%%%%%%%%%%%%%%%%%%%%%%%%%

The analysis yields the following key results:

\begin{enumerate}
    \item \textbf{Result 1:} [Primary finding with quantitative support]
    \item \textbf{Result 2:} [Secondary finding with implications]
    \item \textbf{Result 3:} [Practical application or boundary condition]
\end{enumerate}


\section{Integration with EEE and FFF}
\label{sec:ggg-integration}
%%%%%%%%%%%%%%%%%%%%%%%%%%%%%%%%%%%%%%%%%%%%%%%%%%%%%%%%%%%%%%%%%%%%%%%%%%%%%%%

\begin{tcolorbox}[colback=yellow!5!white, colframe=yellow!75!black,
    title=Integration: METHOD-CONFIG in the Workflow]

\textbf{Position in EEE Workflow:}
\begin{enumerate}
\item Entry Point (Step 1)
\item Scope Definition (Step 2)
\item Context Specification (Step 3)
\item \textbf{→ CONFIGURATOR (GGG) generates initial model}
\item Variable Selection (Step 4) --- refine CONFIG output
\item Functional Form (Step 5) --- refine CONFIG output
\item Parameter Estimation (Step 6)
\item Predictions (Step 7)
\item Register Model in FFF (Step 8)
\end{enumerate}

\textbf{Data Flow:}
\begin{verbatim}
User Input → GGG CONFIG → Initial Model → EEE Refinement → FFF Storage
\end{verbatim}

\textbf{Source Tracking:}
Models derived from CONFIG are tagged in FFF registry as:
\begin{verbatim}
parameters_theta:
  sources:
    C_F: "GGG-CONFIG Table 2, Domain=Vorsorge, L=0"
    gamma_FxE: "GGG-CONFIG Table 3, Domain=Vorsorge"
\end{verbatim}
\end{tcolorbox}


% =============================================================================
%                    PART 3: APPLICATION
% =============================================================================

%%%%%%%%%%%%%%%%%%%%%%%%%%%%%%%%%%%%%%%%%%%%%%%%%%%%%%%%%%%%%%%%%%%%%%%%%%%%%%%
\section{Worked Example}
\label{sec:ggg-example}
%%%%%%%%%%%%%%%%%%%%%%%%%%%%%%%%%%%%%%%%%%%%%%%%%%%%%%%%%%%%%%%%%%%%%%%%%%%%%%%

\subsection{Example: Configuring a Health Adherence Model for Austria}

\paragraph{User Input:}
\begin{itemize}
\item Level: L=0 (Individual)
\item Geography: Austria (AT)
\item Domain: Gesundheit
\item Scope Type: Process (adherence journey)
\end{itemize}

\paragraph{Step 1: $\Psi$ Lookup (Table \ref{tab:ggg-psi-mapping})}
\begin{verbatim}
psi_trust: 0.70
psi_digital: 0.80
psi_autonomy: 0.65
psi_risk: 0.50
psi_regulatory: "ASVG"
\end{verbatim}

\paragraph{Step 2: $C$ Lookup (Table \ref{tab:ggg-c-mapping})}
\begin{verbatim}
C_F: 0.15, C_E: 0.30, C_P: 0.35, C_S: 0.15, C_D: 0.03, C_X: 0.02
\end{verbatim}

\paragraph{Step 3: $\gamma$ Lookup (Table \ref{tab:ggg-gamma-mapping})}
\begin{verbatim}
gamma_PxE: 0.40  (health-wellbeing)
gamma_temporal: 0.35  (present bias)
gamma_PxS: 0.15  (health-social)
\end{verbatim}

\paragraph{Step 4: Form Lookup (Table \ref{tab:ggg-form-mapping})}
\begin{verbatim}
Form: Dynamic/BCJ
dS/dt = (1/tau) * [S*(A, WAX, Psi) - S(t)] + epsilon
\end{verbatim}

\paragraph{Generated Model:}
A BCJ-based adherence model with:
\begin{itemize}
\item Strong Physical-Emotional complementarity ($\gamma_{P \times E} = 0.40$)
\item High present bias ($\gamma_{temporal} = 0.35$)
\item Moderate institutional trust ($\Psi_{trust} = 0.70$)
\item Five-stage journey: Pre-contemplation → Maintenance
\end{itemize}


%%%%%%%%%%%%%%%%%%%%%%%%%%%%%%%%%%%%%%%%%%%%%%%%%%%%%%%%%%%%%%%%%%%%%%%%%%%%%%%
\section{Implications for Practice}
\label{sec:ggg-practice}
%%%%%%%%%%%%%%%%%%%%%%%%%%%%%%%%%%%%%%%%%%%%%%%%%%%%%%%%%%%%%%%%%%%%%%%%%%%%%%%

\begin{tcolorbox}[colback=green!5!white, colframe=green!75!black,
    title=Practical Implications]
\begin{enumerate}
\item \textbf{Use CONFIG as starting point:} Never design from blank slate.
    The configurator provides evidence-based defaults.
\item \textbf{Always refine:} CONFIG output is a \textit{draft}, not final.
    Domain expertise should adjust parameters.
\item \textbf{Document deviations:} When you change CONFIG defaults, document
    why. This builds organizational knowledge.
\item \textbf{Update mapping tables:} When validated models outperform CONFIG
    defaults, update the tables. This is organizational learning.
\item \textbf{Respect the ranges:} $\gamma$ values are given as ranges (e.g., 0.20--0.40).
    Choose within range based on context specifics.
\end{enumerate}
\end{tcolorbox}


% =============================================================================
%                    PART 4: BACK MATTER
% =============================================================================

%%%%%%%%%%%%%%%%%%%%%%%%%%%%%%%%%%%%%%%%%%%%%%%%%%%%%%%%%%%%%%%%%%%%%%%%%%%%%%%
\section{Summary}
\label{sec:ggg-summary}
%%%%%%%%%%%%%%%%%%%%%%%%%%%%%%%%%%%%%%%%%%%%%%%%%%%%%%%%%%%%%%%%%%%%%%%%%%%%%%%

\begin{tcolorbox}[colback=green!10!white, colframe=green!75!black,
    title=\textbf{Appendix UNMAPPED_CFG: Summary}]

\textbf{Core Question:} How do we generate model specifications from high-level choices?

\textbf{Main Contribution:}
\begin{itemize}[nosep]
    \item Defines 4 configuration dimensions (Level, Geography, Domain, Scope)
    \item Provides mapping tables for $\Psi$, $C$, $\gamma$, functional form
    \item Includes 8 DACH reference configurations
    \item Integrates with EEE workflow and FFF registry
\end{itemize}

\textbf{Key Tables:}
\begin{itemize}[nosep]
    \item Table \ref{tab:ggg-psi-mapping}: Level × Geography → $\Psi$
    \item Table \ref{tab:ggg-c-mapping}: Domain × Level → $C$
    \item Table \ref{tab:ggg-gamma-mapping}: Domain → $\gamma$
    \item Table \ref{tab:ggg-form-mapping}: Scope Type → Form
\end{itemize}

\textbf{Integration:} GGG provides input generation for EEE workflow;
generated models are stored in FFF registry.

\end{tcolorbox}


%%%%%%%%%%%%%%%%%%%%%%%%%%%%%%%%%%%%%%%%%%%%%%%%%%%%%%%%%%%%%%%%%%%%%%%%%%%%%%%
\section{Glossary of Symbols}
\label{sec:ggg-glossary}
%%%%%%%%%%%%%%%%%%%%%%%%%%%%%%%%%%%%%%%%%%%%%%%%%%%%%%%%%%%%%%%%%%%%%%%%%%%%%%%

\begin{tcolorbox}[colback=blue!5!white, colframe=blue!75!black]
\textbf{Central Reference:} For complete symbol definitions, see
\textbf{Appendix UNMAPPED_GLS} (\textit{Glossary and Notation}, \S\ref{app:glossary}).

This section lists only symbols introduced or specialized in this appendix.
\end{tcolorbox}

\vspace{0.5em}

\begin{table}[htbp]
\centering
\caption{Symbols Introduced in Appendix UNMAPPED_CFG}
\label{tab:ggg-symbols}
\renewcommand{\arraystretch}{1.3}
\begin{tabular}{@{}clp{6cm}c@{}}
\toprule
\textbf{Symbol} & \textbf{Name} & \textbf{Definition} & \textbf{Also in G?} \\
\midrule
$\Psi_{geo}$ & Geographic Context & Country/region context dimension & NEW \\
$\Psi_{trust}$ & Institutional Trust & Trust in institutions [0,1] & NEW \\
$\Psi_{digital}$ & Digital Adoption & Technology adoption level [0,1] & NEW \\
DACH-$nn$ & Reference Config & Pre-validated configuration ID & NEW \\
\bottomrule
\end{tabular}
\end{table}


%%%%%%%%%%%%%%%%%%%%%%%%%%%%%%%%%%%%%%%%%%%%%%%%%%%%%%%%%%%%%%%%%%%%%%%%%%%%%%%

%%%%%%%%%%%%%%%%%%%%%%%%%%%%%%%%%%%%%%%%%%%%%%%%%%%%%%%%%%%%%%%%%%%%%%%%%%%%%%%
\section{Critical Foundations and Objections}
\label{sec:cfg-foundations}
%%%%%%%%%%%%%%%%%%%%%%%%%%%%%%%%%%%%%%%%%%%%%%%%%%%%%%%%%%%%%%%%%%%%%%%%%%%%%%%

\subsection{Objection 1: Scope Limitations}

\textbf{Objection:} The framework presented may have limited applicability
outside the specific domain contexts discussed.

\textbf{Response:} We acknowledge domain-specificity. The framework provides
general principles that should be calibrated to specific contexts. Cross-domain
validation remains an open research question.

\subsection{Objection 2: Measurement Challenges}

\textbf{Objection:} Key constructs may be difficult to measure reliably in
practice.

\textbf{Response:} We recommend using validated scales and instruments where
available, and developing domain-specific proxies where necessary. Sensitivity
analysis should assess robustness to measurement uncertainty.


\section{Open Issues and Research Agenda}
\label{sec:ggg-open}
%%%%%%%%%%%%%%%%%%%%%%%%%%%%%%%%%%%%%%%%%%%%%%%%%%%%%%%%%%%%%%%%%%%%%%%%%%%%%%%

\begin{table}[htbp]
\centering
\caption{Research Roadmap for METHOD-CONFIG}
\label{tab:ggg-roadmap}
\renewcommand{\arraystretch}{1.3}
\begin{tabular}{@{}clcl@{}}
\toprule
\textbf{\#} & \textbf{Issue} & \textbf{Priority} & \textbf{Status} \\
\midrule
1 & Empirical validation of mapping values & HIGH & Planned \\
2 & Expansion to non-DACH geographies & MEDIUM & Future \\
3 & Machine learning for mapping optimization & LOW & Research \\
4 & Cross-level interaction mappings & MEDIUM & Open \\
5 & Temporal dynamics in $\Psi$ mappings & LOW & Exploratory \\
\bottomrule
\end{tabular}
\end{table}


%%%%%%%%%%%%%%%%%%%%%%%%%%%%%%%%%%%%%%%%%%%%%%%%%%%%%%%%%%%%%%%%%%%%%%%%%%%%%%%
\section{References}
\label{sec:ggg-references}
%%%%%%%%%%%%%%%%%%%%%%%%%%%%%%%%%%%%%%%%%%%%%%%%%%%%%%%%%%%%%%%%%%%%%%%%%%%%%%%

\begin{tcolorbox}[colback=blue!5!white, colframe=blue!75!black]
\textbf{Master Bibliography:} For complete reference details, see
\texttt{bcm2\_master\_references.bib} in the repository.

This section lists appendix-specific references and data sources.
\end{tcolorbox}

\vspace{0.5em}

\subsection{Key References for This Appendix}

\begin{itemize}[nosep]
    \item \textbf{DACH Cultural Dimensions:} Hofstede (2001), GLOBE Study (House et al., 2004)
    \item \textbf{Pension Systems:} OECD Pensions at a Glance (2023)
    \item \textbf{Health Behavior:} Schwarzer (2008) HAPA Model
    \item \textbf{Configuration Design:} Product configuration literature
\end{itemize}

\nocite{bcm_master}


%%%%%%%%%%%%%%%%%%%%%%%%%%%%%%%%%%%%%%%%%%%%%%%%%%%%%%%%%%%%%%%%%%%%%%%%%%%%%%%
% CROSS-REFERENCE FOOTER (Required)
%%%%%%%%%%%%%%%%%%%%%%%%%%%%%%%%%%%%%%%%%%%%%%%%%%%%%%%%%%%%%%%%%%%%%%%%%%%%%%%

\vspace{1em}
\hrule
\vspace{0.5em}

\noindent\textit{Cross-references:}
Appendix UNMAPPED_GLS (Central Glossary),
Appendix UNMAPPED_DSN (METHOD-DESIGN),
Appendix UNMAPPED_REG (METHOD-REGISTRY),
Appendix UNMAPPED_WHO (CORE-WHO),
Appendix CTW (CORE-WHEN),
Appendix UNMAPPED_WAT (CORE-WHAT).

\noindent\textit{Master Bibliography:} \texttt{bcm2\_master\_references.bib}


% =============================================================================
% END OF APPENDIX GGG
% =============================================================================
